\section{Limitations and Scope}
\label{sec:limitations}

\systemname\ has five principal limitations. First, it keeps kernel-ready high- and low-precision expert representations in host memory. This removes online repacking from the critical path but increases pinned-DRAM demand and may be unsuitable for edge hosts with limited memory. Disk-backed preparation would reduce DRAM use at the cost of another latency tier.

Second, the current policy optimizes aggregate routing utility. It does not provide per-tenant quality guarantees or fairness when concurrent requests have different domains. Per-request precision decisions would also make batching less regular. Extending the objective with tenant weights or service-level constraints is therefore a policy problem that the present mechanism does not solve.

Third, our evaluation covers one 48\,GB GPU and three model families. Transfer/compute overlap and the feasible high-precision capacity depend on the GPU, PCIe generation, sequence length, and KV-cache policy, so the operating points are platform measurements. The pinned MoE-Infinity open-source runtime supports Qwen3-30B but not the other two checkpoints, and it differs from the implementation in its paper. The design uses two precomputed tiers on one GPU and a fixed task order.

Fourth, the prototype favors correctness and inspectability over kernel maturity. Handle-aware dispatch currently uses eager expert execution for Qwen-family models because the upstream grouped-GEMM shortcut does not accept per-expert versioned handles. This isolates the systems effect under study but leaves integration with a fused handle-aware grouped kernel to future work. Dual-tier packing and pinned-host-cache construction are untimed startup costs; production deployments would need to amortize or persist this preparation.

Finally, the hard invariant applies to \systemname-owned expert pools, staging blocks, and declared conversion workspace. End-to-end device-memory safety is conditional on the serving configuration keeping KV-cache, activation, and framework allocations within their startup reservations; the prototype does not replace or police the framework allocator. A separate CUDA stream removes explicit migration synchronization from the forward path but cannot eliminate resource contention: copies and repacking kernels may still share memory-system bandwidth with inference. Our latency scope is an isolated model and excludes request queueing, tokenization, network transport, and serialization. We therefore compare model-level end-to-end and tail latency, not only scheduler CPU time.
